\documentclass[10pt]{article}
\providecommand{\FIGDIR}{fig}
\newcommand{\DIGESTDATE}{9 September 2026}
% ===================== PM Research Watch - shared preamble =====================
\usepackage[T1]{fontenc}
\usepackage[utf8]{inputenc}
\usepackage{amsmath,amssymb}
\usepackage{tgpagella}
\usepackage{tgheros}
\renewcommand{\sfdefault}{qhv}
\renewcommand{\ttdefault}{lmtt}
\usepackage{microtype}
\usepackage[english]{babel}

\usepackage[a4paper,top=17mm,bottom=17mm,left=17mm,right=17mm,
            headsep=6mm,footskip=10mm]{geometry}
\usepackage{graphicx}
\usepackage{xcolor}
\usepackage{colortbl}
\usepackage{tikz}
\usetikzlibrary{calc}
\usepackage[most]{tcolorbox}
\usepackage{titlesec}
\usepackage{fancyhdr}
\usepackage{booktabs}
\usepackage{array}
\usepackage{longtable}
\usepackage{enumitem}
\usepackage{multicol}
\usepackage{ragged2e}
\usepackage[hidelinks]{hyperref}

% ------------------------------------------------------------------ palette
\definecolor{Ink}    {HTML}{14262E}
\definecolor{Deep}   {HTML}{0F4C5C}
\definecolor{Teal}   {HTML}{1F7A8C}
\definecolor{Sky}    {HTML}{4FA3C4}
\definecolor{Sage}   {HTML}{6FA287}
\definecolor{Amber}  {HTML}{E3A008}
\definecolor{Clay}   {HTML}{B5643C}
\definecolor{Coral}  {HTML}{D1495B}
\definecolor{Violet} {HTML}{6A5B9E}
\definecolor{Slate}  {HTML}{4A6572}
\definecolor{Paper}  {HTML}{FBF9F5}
\definecolor{Mist}   {HTML}{EEEAE2}
\definecolor{Rule}   {HTML}{D9D3C9}

\pagecolor{Paper}
\color{Ink}

\hypersetup{colorlinks=true, linkcolor=Deep, urlcolor=Teal, citecolor=Teal,
            pdftitle={PM Research Watch}, pdfauthor={Automated daily digest}}

% ------------------------------------------------------------------ headings
\titleformat{\section}
  {\sffamily\bfseries\large\color{Deep}}{}{0pt}
  {}[{\vspace{1.2mm}\color{Amber}\titlerule[1.1pt]}]
\titlespacing*{\section}{0pt}{5.5mm}{2.8mm}

\titleformat{\subsection}
  {\sffamily\bfseries\normalsize\color{Teal}}{}{0pt}{}
\titlespacing*{\subsection}{0pt}{3.6mm}{1.4mm}

\setlist[itemize]{leftmargin=4.2mm, itemsep=1.1mm, topsep=1.3mm, parsep=0pt,
                  label=\textcolor{Teal}{\footnotesize$\blacktriangleright$}}

% ------------------------------------------------------------------ header/footer
\pagestyle{fancy}
\fancyhf{}
\renewcommand{\headrulewidth}{0pt}
\renewcommand{\footrulewidth}{0pt}
\fancyhead[L]{\sffamily\scriptsize\color{Slate}PM RESEARCH WATCH\;\textcolor{Rule}{\textbar}\;\DIGESTDATE}
\fancyhead[R]{\sffamily\scriptsize\color{Slate}Particulate matter monitoring \& health effects}
\fancyfoot[C]{\sffamily\scriptsize\color{Slate}\thepage}

% ------------------------------------------------------------------ components
% metric tile
\newcommand{\metric}[3]{%
  \begin{tcolorbox}[enhanced, colback=white, colframe=Rule, boxrule=0.5pt,
    arc=1.6mm, left=2.2mm, right=2.2mm, top=1.8mm, bottom=1.8mm,
    borderline west={1.7mm}{0pt}{#3}, nobeforeafter, valign=center,
    height=17mm, width=\linewidth]
    \RaggedRight
    {\sffamily\bfseries\fontsize{17}{17}\selectfont\color{#3}#1}\\[0.5mm]
    {\sffamily\fontsize{6.5}{7.4}\selectfont\color{Slate}#2}
  \end{tcolorbox}}

% pull-quote / key insight box
\newtcolorbox{keybox}[1][]{enhanced, colback=Mist, colframe=Mist,
  arc=1.6mm, left=3.2mm, right=3.2mm, top=2.4mm, bottom=2.4mm,
  borderline west={1.7mm}{0pt}{Amber}, #1}

% subtopic band
\newcommand{\band}[2]{%
  \vspace{3.2mm}\noindent
  \begin{tcolorbox}[enhanced, colback=#2, colframe=#2, arc=1.2mm,
    left=2.6mm, right=2.6mm, top=1.1mm, bottom=1.1mm, boxrule=0pt,
    width=\linewidth, nobeforeafter]
    {\sffamily\bfseries\small\color{white}#1}
  \end{tcolorbox}\vspace{1.6mm}}

% one paper entry: {short cite}{title}{journal}{doi}{body}
\newcommand{\paperentry}[5]{%
  \par\vspace{1.9mm}\noindent
  {\sffamily\bfseries\footnotesize\color{Ink}#1}\;%
  {\sffamily\fontsize{7}{8}\selectfont\color{Slate}#3}\\[0.4mm]
  {\footnotesize\itshape\color{Deep}#2}\\[0.6mm]
  {\footnotesize #4}\\[0.3mm]
  {\fontsize{6.6}{7.6}\selectfont\color{Slate}\texttt{doi:}\,\href{https://doi.org/#5}{\texttt{\detokenize{#5}}}}
  \par}

% caption style
\newcommand{\figcap}[1]{\par\vspace{1.0mm}
  {\fontsize{7.2}{8.6}\selectfont\color{Slate}\textsf{\textbf{Fig.}}\; #1\par}\vspace{1.5mm}}

\providecommand{\permil}{\textperthousand}

\setlength{\parindent}{0pt}
\setlength{\parskip}{0pt}
\renewcommand{\arraystretch}{1.18}


\begin{document}

% ==================================================================== MASTHEAD
\thispagestyle{empty}
\begin{tikzpicture}[remember picture, overlay]
  \fill[Deep] (current page.north west) rectangle
      ($(current page.north east)+(0,-67mm)$);
  \fill[Amber] ($(current page.north west)+(0,-67mm)$) rectangle
      ($(current page.north east)+(0,-68.4mm)$);
  % subtle particle field
  \foreach \x/\y/\r in {22/12/0.9, 41/26/1.5, 63/9/0.7, 88/31/1.1, 112/17/1.9,
                        138/29/0.8, 159/13/1.3, 178/54/1.0, 196/21/1.6,
                        30/44/0.6, 74/57/1.2, 125/48/0.9, 168/6/0.8, 191/59/0.7,
                        52/18/1.0, 100/60/0.6, 148/20/1.1, 66/49/1.4, 205/45/0.8}
    \fill[Sky, opacity=0.30] ($(current page.north west)+(\x mm,-\y mm)$) circle (\r mm);
\end{tikzpicture}

\vspace*{4mm}
{\sffamily\fontsize{7.6}{9}\selectfont\color{Sky}%
 AUTOMATED DAILY LITERATURE DIGEST\;\textbullet\;PhD RESEARCH WATCH\par}
\vspace{2.2mm}
{\sffamily\bfseries\fontsize{27}{29}\selectfont\color{white}%
 Particulate Matter\par}
{\sffamily\fontsize{27}{29}\selectfont\color{Sky}%
 Monitoring \& Health Effects\par}
\vspace{5.0mm}
{\sffamily\fontsize{8.4}{10}\selectfont\color{white}%
 Issue 2026-09-09\;\;\textcolor{Amber}{\textbar}\;\;Entry window 9 September 2026\;\;%
 \textcolor{Amber}{\textbar}\;\;26 records screened in scope\par}
\vspace{18mm}

% -------------------------------------------------------------- metric strip
\noindent
\begin{minipage}[t]{0.190\linewidth}\metric{26}{RECORDS IN SCOPE\\(71 EXCLUDED)}{Deep}\end{minipage}\hfill
\begin{minipage}[t]{0.190\linewidth}\metric{9}{SUBTOPIC CLUSTERS\\REPRESENTED}{Teal}\end{minipage}\hfill
\begin{minipage}[t]{0.190\linewidth}\metric{10}{EFFECT ESTIMATES\\EXTRACTED}{Sky}\end{minipage}\hfill
\begin{minipage}[t]{0.190\linewidth}\metric{7}{SENSING \&\\INSTRUMENTATION}{Amber}\end{minipage}\hfill
\begin{minipage}[t]{0.190\linewidth}\metric{4}{METADATA ONLY\\(NO ABSTRACT)}{Clay}\end{minipage}

\vspace{6mm}

% -------------------------------------------------------------- signal of day
\section*{\sffamily\bfseries\color{Deep}Signal of the day}
\vspace{-1.5mm}{\color{Amber}\rule{\linewidth}{1.1pt}}\vspace{2.5mm}

\begin{keybox}
{\sffamily\bfseries\small\color{Deep}The PM$_{2.5}$ component that carries the health
signal is the one the field's standard instrument reads wrong.}\\[1.4mm]
\footnotesize
Two records that arrived on the same day, from literatures that do not cite each other.
Massabò et al.\ (\textit{Atmosphere}) ran an AE33 aethalometer against two independent
offline filter techniques --- MWAA and BLAnCA --- at an urban coastal site in Genoa. The
three agree on \emph{shape} (\textbf{$R^2$ 0.93--0.99} across UV, visible and near-IR) and
disagree on \emph{level}: AE33 reports absorption \textbf{12--22\% high}, wavelength
dependent, which points at the default multiple-scattering correction rather than at
noise. Meanwhile Yu et al.\ (\textit{JAHA}), in a time-stratified case-crossover on nine
years of Hangzhou acute aortic dissections, found that an IQR rise in black carbon carried
a \textbf{42\% increase} in type-A dissection odds (OR 1.42, 95\% CI 1.15--1.75), ahead of
organic matter (1.35) and sulfate (1.36) --- and that BC alone took
\textbf{over half the weighted-quantile-sum joint weight}. So the component-resolved
epidemiology is converging on BC exactly where the routine BC measurement carries a
double-digit multiplicative bias that varies by site and aerosol population. A 15\% error
in the exposure is not a rounding detail when the effect being estimated is a 42\% odds
increase per IQR. The third record closes the loop: the RI-URBANS synthesis
(\textit{Atmos.\ Meas.\ Tech.}) reports, across a continent of low-cost-sensor, mobile and
citizen-science pilots, that the binding constraint on usable urban maps was never sensor
accuracy --- it was \emph{harmonisation} between methods. For distributed sensing the
implication is concrete and unwelcome: a dense BC network calibrated against a
site-specific aethalometer offset is not comparable to anyone else's dense BC network.
\end{keybox}

\vspace{3mm}

% -------------------------------------------------------------- exec brief
\section*{\sffamily\bfseries\color{Deep}What changed today}
\vspace{-1.5mm}{\color{Amber}\rule{\linewidth}{1.1pt}}\vspace{2.5mm}

\footnotesize
\begin{itemize}
\item \textbf{Personal and mobile monitoring produced the two strongest exposure records
of the day, and they disagree about what resolution buys you.} MobiliSense
(\textit{preprint}) put two ambulatory monitors, a GPS and an accelerometer on 199 Paris
adults and took 2{,}504 spirometry tests: BC one hour before testing cut FEV$_1$ by
\textbf{0.016\,L} (95\% CI $-$0.024 to $-$0.008) per 1\,$\mu$g/m$^3$, and PM$_{2.5}$ at
15--60 minutes cut FEV$_1$/FVC by 0.5--0.7 percentage points --- effects that a daily
residential average cannot resolve at all. The Hong Kong GEMA study
(\textit{Environ Res}) went further, 1.7 million personal PM$_{2.5}$ measurements on 800
adults across eight activity-travel contexts, and found that context-specific
associations \emph{reversed sign} while the pooled estimate showed nothing sustained.
Higher resolution is not automatically a stronger signal; it can be a decomposition of a
null into opposing parts.

\item \textbf{Three clean nulls, all of them worth more than a marginal positive.}
A 24-cohort, 126-million-participant meta-analysis (\textit{Environ Pollut}) found no sex
modification of the PM$_{2.5}$--Alzheimer's association (RRR 0.963, 0.919--1.009) or of
vascular dementia (1.003, 0.954--1.054), with GRADE moderate certainty, and states
plainly that sex-specific policy is not supported. A matched case-control with a
distributed-lag model found the green-space--ASD protective window is real in mid-to-late
pregnancy and that PM$_{2.5}$ \emph{neither mediated nor modified} it --- closing the
obvious confounding pathway rather than leaving it open. And the smoke-free-home cluster
RCT missed its primary outcome outright.

\item \textbf{Non-exhaust and non-filterable sources both surfaced, from opposite
directions.} Tire-wear particles inhibited shoot and root growth and raised root lipid
peroxidation in a roadside grass (\textit{Curr Biol}); condensable PM --- which leaves the
stack as vapour and forms particles downstream, so filterable-PM controls miss it entirely
--- was removed at 29.6--65.3\% by a string-of-beads liquid at under a tenth the energy of
condensation (\textit{J Hazard Mater}). Both are sources that mass-based optical monitors
size or miss badly.

\item \textbf{The abstract-deposit gap is now a structural property of a one-day window,
not an incident.} Four Elsevier records (\textit{Environ Int}, \textit{Atmos Environ},
\textit{Atmos Pollut Res}, \textit{Eur J Intern Med}) had no abstract in Crossref per-DOI,
Europe PMC or PubMed on the entry date. They ship as \textsc{metadata only}, asserting no
finding. This is the second consecutive issue to hit it; see the provenance note.

\item \textbf{The PubMed connector added net-new records for the first time.} Its health
axis returned 16 PMIDs against the harvester's own E-utilities leg's 11, nine of them
unseen, three of which ship --- an analytical method for atmospheric alkylamines at
0.1--1.1\,pg/m$^3$ in PM, the condensable-PM removal study, and a cardiorespiratory
review. On 8 September the connector was pure confirmation. It is not redundant.
\end{itemize}

\clearpage

% ==================================================================== ANALYTICS
\section*{\sffamily\bfseries\color{Deep}Corpus analytics}
\vspace{-1.5mm}{\color{Amber}\rule{\linewidth}{1.1pt}}\vspace{3mm}

\noindent\includegraphics[width=\linewidth]{\FIGDIR/f1_subtopics.png}
\figcap{Subtopic assignment of the 26 in-scope records; single-label by dominant
contribution. \textit{Sensing, forecasting \& instrumentation} is the largest cluster at
7, followed by \textit{Exposure assessment \& modelling} at 4 --- together 42\% of the
issue, an instrumentation-weighted day. The four clinical clusters
(cardiovascular/metabolic, neuro, respiratory, reproductive) carry two records each.
\textit{Other clinical endpoints} is the only empty cluster: \textbf{9 of 10 represented}
ties 4 September for the broadest coverage of the past week, against 6 clusters on
8 September.}

\vspace{2mm}
\noindent\includegraphics[width=\linewidth]{\FIGDIR/f2_design_endpoint.png}
\figcap{Left --- study architecture. \textit{Modelling / inventory} and
\textit{Metadata only} tie at 4 each; the three observational headings
(cross-sectional 3, acute 2, cohort 1) sum to \textbf{6}, with one further
\textit{Ecological} record. Two trial records is unusual for a single day, and neither
reports a positive primary result: one is a protocol, the other missed its endpoint.
Right --- health endpoint. \textbf{Eleven of 26 records carry no human health endpoint at
all} --- the monitoring, instrument, emissions and modelling records --- against 4
respiratory, 3 cardiovascular and 2 each for mental health and neuro/cognitive. One record
is ecotoxicological (plants, not people) and is labelled as such rather than folded into a
human category.}

\vspace{2mm}
\noindent\includegraphics[width=\linewidth]{\FIGDIR/f3_metric_geography.png}
\figcap{Left --- particle metric. The \textit{PM (unspeciated) / emissions} bar dominates
at 13 because the day is instrumentation- and source-heavy; PM$_{2.5}$ alone is the
exposure in 7 and joint PM$_{2.5}$/PM$_{10}$ in 3. Only \textbf{one} record resolves
composition and one resolves optical properties --- which is the gap the Signal of the day
turns on. Right --- geographic provenance. Europe and the
\textit{Global / multi-region} fallback tie at 7, but the fallback is not what it looks
like: 3 are genuinely multi-region reviews and \textbf{4 are records whose study setting
could not be established} because they have no abstract. China contributes 6. South Asia,
Oceania, MENA and East Asia (ex-China) contribute one apiece; sub-Saharan Africa
contributes none, against 3 on 8 September.}

\vspace{2mm}
\noindent\includegraphics[width=\linewidth]{\FIGDIR/f6_lifecourse.png}
\figcap{Life-course windows addressed; counts are non-exclusive and only the 11 records
with an identifiable human window are placed. Working-age adulthood dominates at 5 --- the
Hong Kong and Paris mobility cohorts, the permanent-supportive-housing trial, the rescue
personnel survey and the CFRP occupational review. Puberty and adolescence is empty, as it
has been for most of this archive; the two in-utero records are the prenatal green-space
window study and the preeclampsia record.}

\vspace{2mm}
\noindent\includegraphics[width=\linewidth]{\FIGDIR/f4_heatmap.png}
\figcap{Subtopic $\times$ study-architecture cross-tabulation. The single densest cell is
\textit{Sensing} $\times$ \textit{Measurement campaign} at 3 (the Genoa intercomparison,
the Baltic ship-plume campaign and the RI-URBANS synthesis); three further cells tie at 2.
The structural feature to note is the same one this archive keeps recording: \textbf{all
seven \textit{Sensing} records carry no health endpoint}, and three of the four
\textit{Exposure} records do not either --- the sole exception in these two rows being the
Hong Kong GEMA study. Only two records in the whole issue put a \emph{measured personal}
particle exposure and a human outcome in the same model, and the other one (MobiliSense)
is filed under \textit{Respiratory}. In the other direction every clinical row is empty of
measurement-campaign work. The two literatures met on one page today mostly because a
digest put them there.}

\vspace{2mm}
\noindent\includegraphics[width=\linewidth]{\FIGDIR/f5_forest.png}
\figcap{All ten quantitative estimates in the corpus on a common logarithmic axis.
\textbf{The contrasts are not mutually comparable and three different quantities are
plotted.} The three aortic-dissection ORs are per interquartile-range increments of
distinct PM$_{2.5}$ components; the pulmonary-hypertension HR is per 1\,SD of a composite
pollution score; the three \textit{RRR} points are \emph{ratios of risk ratios} --- female
versus male --- so for those the null line means \textbf{no sex difference}, not no
effect, and reading them as small associations would be wrong. The two lowest-precision
entries are interventions, not exposures: the smoke-free-home trial's primary outcome
(OR 1.45, 0.69--3.07) crosses the null and its any-adoption secondary (3.83, 2.63--5.57)
does not. The passive-smoking estimate is the univariate one, shown because the paper's
own multivariate interval (1.959--3.144 around OR 1.990) is internally inconsistent. The
trial's CO-verified abstinence OR of 6.94 (1.69--28.45) is omitted: it rests on 14 events
total and would have compressed every other estimate on the axis.}

\clearpage

% ==================================================================== DIGEST
\section*{\sffamily\bfseries\color{Deep}Paper-level digest}
\vspace{-1.5mm}{\color{Amber}\rule{\linewidth}{1.1pt}}\vspace{1mm}

{\fontsize{7.4}{9}\selectfont\color{Slate}Ordered by cluster. Each entry gives the
methodological core and the finding that matters, not an abstract paraphrase. Entries
marked \textsc{metadata only} had no abstract in existence on the entry date; for these
the entry states what the record establishes and what it does not, and asserts no
finding.\par}

\band{Sensing, forecasting \& instrumentation\hfill 7 records}{Amber}

\paperentry{Genoa absorption intercomparison 2026}
{AE33 aethalometer against two filter-based absorption techniques in a coastal city}
{Atmosphere}
{Multi-wavelength field intercomparison at the Multedo urban coastal site, Genoa, where the
aerosol mixes traffic, shipping, industry and marine air. The AE33 was run against the
Multi-Wavelength Absorption Analyzer and the Broadband Light Analyzer of Complex Aerosols
across UV, visible and near-IR. Correlations are excellent --- \textbf{$R^2$ 0.93--0.99},
tightest between the two offline methods --- but AE33 reports absorption coefficients
\textbf{12--22\% higher}, with the offset varying by wavelength. That pattern implicates
the instrument's default multiple-scattering correction, which is tuned to a generic
aerosol and not to this one. Absorption Ångström exponents from all three sit near unity,
consistent with dominant primary combustion. The limitation is that one site and one
campaign cannot fix the correction factor generally --- which is precisely the point:
a site-specific bias is not a constant you can publish once.}
{10.3390/atmos17090883}

\paperentry{RI-URBANS urban AQ mapping 2026}
{What a continent of low-cost-sensor and citizen-science air-quality pilots actually taught}
{Atmos Meas Tech}
{Synthesis of the RI-URBANS mapping and hotspot-identification pilots across European
cities, covering stationary low-cost sensor networks, mobile monitoring platforms and
citizen-collected data alongside reference networks and modelling. The conclusion the
programme reaches is organisational rather than technical: the limiting factor on
producing comparable urban maps was \textbf{harmonisation between methods}, not the
accuracy of any one method, and harmonisation was one of the project's stated primary
aims for that reason. Two operational findings are worth carrying into any deployment ---
campaign design and sustained, direct communication with citizen participants determined
data quality as much as hardware did. As a programme-level synthesis it reports no single
effect estimate and cannot be appraised like a primary study; its value is the negative
result about where the effort goes.}
{10.5194/amt-19-5729-2026}

\paperentry{Alkylamine IC-MS/MS method 2026}
{Speciated atmospheric alkylamines at picogram levels in gas and particle phase}
{J Chromatogr A}
{Alkylamines drive new particle formation and early aerosol growth, and are analytically
awkward --- strongly polar, present at trace level, prone to adsorption and decomposition.
This method pairs alkaline purging with acidic absorption for selective enrichment ahead of
ion chromatography--tandem MS. Method detection limits are \textbf{9--94\,pg\,m$^{-3}$
gaseous and 0.1--1.1\,pg\,m$^{-3}$ in PM}, a substantial gain on existing offline methods,
and the approach resolves the dimethylamine/monoethylamine isobaric pair and quantifies
tertiary amines that derivatisation-based methods structurally cannot. Field-applied at
septic tanks and traffic tunnels in Dalian, it recovered a broader speciation profile than
previous work and seasonal variation in each species. It is an offline enrichment method,
so time resolution is whatever the sampling interval is --- this does not compete with
online amine measurement for process studies, it competes on speciation and sensitivity.}
{10.1016/j.chroma.2026.467416}

\paperentry{Baltic ship plume OFR aging 2026}
{Fresh and photochemically aged emission factors from individual ships, measured at sea}
{Environ Sci Technol Lett}
{An oxidation flow reactor operated from a moving research vessel, sampling identified
ocean-going ships at \textbf{0.6--1.6\,km} standoff in the Baltic, plus the research
vessel's own plume as a control. Gas-phase and potential-aerosol-mass emission factors
were measured fresh and after photochemical aging. \textbf{Particle mass rose after aging
for every vessel sampled}, and secondary aerosol formation potential scaled with SO$_2$
emission --- i.e.\ with fuel sulfur content, consistent with sulfate-driven particle
formation. This is a pilot with a small vessel sample and no size-resolved chemistry
reported in the abstract, so the emission factors are demonstrative rather than
inventory-grade. The contribution is the framework: stack-level aging experiments and
fixed ambient monitoring have both existed for years, and this joins them on a platform
that can follow an individual, identifiable source.}
{10.1021/acs.estlett.6c00695}

\paperentry{India winter PM2.5 seasonal forecast 2026}
{Winter PM$_{2.5}$ severity over northern India is predictable a season ahead}
{Sci Adv}
{Over a decade of quality-controlled Indian government reference PM$_{2.5}$ monitoring,
used to establish that \textbf{roughly 70\% of interannual variance} in northern Indian
winter PM$_{2.5}$ is explained by western disturbances --- Mediterranean-origin
extratropical storms --- acting through ventilation and wet removal. The storms in turn
track to the position and strength of the subtropical jet, which is traceable to North
Atlantic and tropical Indian Ocean sea-surface-temperature anomalies in the preceding
autumn. Using only preseason SST, the authors predict winter WD precipitation and PM$_{2.5}$
severity one season ahead at \textbf{$r^2$ 0.52--0.69}. The honest caveat is that this is
meteorological predictability, not emissions predictability: it forecasts how well a given
emission load will be ventilated, and a decade of records is a short sample for an
interannual teleconnection. For network planning it is nonetheless actionable --- it says
which winters to instrument heavily.}
{10.1126/sciadv.aee4549}

\paperentry{US precipitation washout DoubleML 2026}
{Observation-derived washout coefficients, and how wildfire smoke breaks them}
{Research Square}
{\textbf{Preprint.} Ten years of hourly EPA observations --- 760 PM$_{2.5}$ and 451
PM$_{10}$ monitors across the contiguous US --- merged with NLDAS-2 meteorology and NOAA
HMS smoke classifications, with a double machine-learning causal framework separating
precipitation from confounding meteorology. Removal rises with rainfall intensity,
particle size and initial concentration, and \textbf{falls as wildfire smoke loading
rises} --- smoke aerosol is less efficiently scavenged, which matters more each year. The
authors fit power-law washout coefficients ($a=0.0404$, $b=0.5951$ for PM$_{2.5}$;
$a=0.0812$, $b=0.4508$ for PM$_{10}$) intended to drop straight into chemical transport
models. The inference threat is the usual one for DoubleML on observational data:
unmeasured confounding is assumed away rather than eliminated, and hourly monitor-scale
precipitation from a reanalysis is a coarse proxy for what fell on the sampler.}
{10.21203/rs.3.rs-10855101/v1}

\paperentry{Open-pit mine dust stacking model 2026}
{\textsc{metadata only} --- multi-point cold-season dust prediction at an open-pit coal mine}
{Atmos Pollut Res}
{\textbf{No abstract existed in Crossref queried per-DOI, in Europe PMC or in PubMed on the
entry date.} What the record establishes: a multi-point dust-concentration prediction task
at an open-pit coal mine during the cold season has been addressed with a stacking ensemble
and an interpretability analysis. What it does not establish: the number and siting of
monitoring points, the prediction horizon, the base learners in the stack, the
interpretability method, any error figure, or the mine's location --- none of which is
inferable from the title and none of which is asserted here. Flagged for retrieval; open-pit
mining dust is one of the few settings where a dense low-cost network has an unambiguous
operational customer.}
{10.1016/j.apr.2026.103196}

\band{Exposure assessment \& modelling\hfill 4 records}{Amber}

\paperentry{Hong Kong GEMA PM2.5 mobility 2026}
{1.7 million personal PM$_{2.5}$ measurements against momentary stress and activity}
{Environ Res}
{Two-day geographic ecological momentary assessment on 800 Hong Kong adults, integrating
\textbf{1.7 million personal PM$_{2.5}$ measurements}, 2.1 million accelerometer records,
4{,}800 momentary assessments and eye-level greenspace derived from 4.3 million
street-view images, summarised over 5- to 600-minute windows and analysed across eight
activity-travel contexts with a DAG-guided multi-relation double-ML framework. Higher
PM$_{2.5}$ tracked more light-intensity and less moderate-to-vigorous activity overall,
and predicted worse \emph{perceived} air quality; perceived pollution in turn predicted
higher stress while measured PM$_{2.5}$ did not. The result that should travel:
\textbf{context-specific estimates ran in opposite directions} while cumulative
associations with stress showed nothing sustained. Two days per participant is thin for a
stress outcome and reverse causation between activity and exposure is unresolved, but the
decomposition itself is the finding --- a pooled null here is an artefact of aggregation.}
{10.1016/j.envres.2026.125659}

\paperentry{NZ tractor fleet PM inventory 2026}
{An uncertainty-propagated PM$_{2.5}$ inventory for an entirely unregulated source class}
{Sci Total Environ}
{Non-road mobile machinery is under-regulated and rising in relative importance as on-road
sources are controlled; New Zealand is an instructive case because it pairs ultra-low-sulphur
diesel and tightening on-road standards with \emph{no} NRMM emission regulation at all. The
inventory joins row-level registration microdata for 43{,}367 tractors, a licensing table
defining 31{,}253 active units, and 29{,}278 dealer sales records (2017--2025), under an
IPCC Tier 2 fuel-anchored framework with EMEP/EEA factors, tier hours calibrated against a
CARB monitoring study, and fuel anchored to a 204\,Ml national consumption target.
10{,}000-iteration Monte Carlo gives \textbf{115\,t PM$_{2.5}$ (95\% UI 73--190)} and
2258\,t NO$_x$. The distributional result is the useful one: \textbf{Gini 0.51 for
PM$_{2.5}$, with the dirtiest 30\% of tractors producing 69\% of it}, and modernisation
running at only $+2.7$ percentage points a year. Activity hours remain the weakest link ---
they are calibrated, not measured.}
{10.1016/j.scitotenv.2026.182302}

\paperentry{Paris Saharan dust heatwave 2026}
{Dust aerosol cools Paris by up to 1\,$^{\circ}$C --- and that does not mean comfort}
{Atmos Chem Phys}
{Three Meso-NH simulations of the 15--19 June 2022 Paris heatwave --- no dust, CAMS
reanalysis dust, and doubled dust --- validated against the PANAME-Urban campaign, which
the model reproduces for aerosol optical depth and incoming solar radiation, with improved
air temperature and boundary-layer height once aerosol is included. The dust plume's
radiative effect cuts surface shortwave by up to \textbf{75\,W\,m$^{-2}$} and air
temperature by up to \textbf{1\,$^{\circ}$C}. The result worth the page is the one that
refuses the easy conclusion: thermal comfort improves by up to 1\,$^{\circ}$C in
\emph{sunlit} urban and suburban areas at 16:00 UTC, but in shaded suburban areas the
increase in diffuse radiation and humidity plus the drop in wind speed cancels the
temperature benefit entirely. A single five-day episode in one city, and comfort indices
are not health outcomes --- but the case for coupling aerosol forecasts to heat-health
warning systems is made on the disaggregation, not on the mean.}
{10.5194/acp-26-12729-2026}

\paperentry{Construction equipment PM-bound EPFR 2026}
{\textsc{metadata only} --- persistent free radicals on PM from non-road construction plant}
{Atmos Environ}
{\textbf{No abstract retrievable on the entry date} from Crossref per-DOI, Europe PMC or
PubMed. What the record establishes: PM-bound environmentally persistent free radicals and
downstream reactive radical species have been measured from non-road construction equipment
under \emph{real-world} rather than dynamometer conditions, which is the harder and rarer
measurement. What it does not establish: radical concentrations or lifetimes, the engine
tiers and duty cycles sampled, the EPR or spin-trapping protocol, or the site --- none is
asserted here. Relevant because EPFR content is a redox-activity metric that mass
concentration does not predict, and construction plant sits inside the urban domain that
dense networks are deployed to characterise.}
{10.1016/j.atmosenv.2026.122353}

\clearpage

\band{Cardiovascular \& metabolic\hfill 2 records}{Teal}

\paperentry{Hangzhou AAD PM components 2026}
{Which PM$_{2.5}$ component triggers acute aortic dissection, and by how much}
{J Am Heart Assoc}
{Time-stratified case-crossover on Hangzhou acute aortic dissections, 2015--2023, with
conditional logistic regression for single components and weighted quantile sum for the
mixture. At \textbf{lag 0--1 days}, interquartile-range increases in black carbon, organic
matter and sulfate were associated with \textbf{42\% (OR 1.42, 1.15--1.75)}, 35\% (1.35,
1.10--1.64) and 36\% (1.36, 1.09--1.69) higher odds of Stanford \emph{type A} dissection;
type B associations were positive but not significant. In the joint model
\textbf{black carbon carried more than half the weight}. Age, smoking, alcohol use and
hypertension modified the associations. The case-crossover design self-controls for
time-invariant confounding, which is its strength; the weaknesses are that component
concentrations are strongly collinear (WQS mitigates but does not resolve this), exposure
is assigned from ambient monitors rather than personally, and a single city's aerosol
composition determines which component appears to dominate.}
{10.1161/jaha.126.049278}

\paperentry{UK PH metabolomic mediation 2026}
{Air pollution, pulmonary hypertension, and how much of it runs through metabolism}
{Ecotoxicol Environ Saf}
{Prospective cohort of \textbf{444{,}346} participants free of pulmonary hypertension at
baseline with NMR metabolomics, 2{,}328 incident cases over a 13.58-year median follow-up.
PM$_{2.5}$, PM$_{10}$, NO$_2$ and NO$_x$ entered individually and as a joint air pollution
score. The score gave \textbf{HR 1.12 per SD (95\% CI 1.08--1.17)}; individual pollutants
ran 1.06--1.93 per 10\,$\mu$g/m$^3$. Elastic-net regression produced a 105-metabolite
signature dominated by lipoprotein measures, fatty acids and amino acids, itself associated
with incident PH (HRs 1.12--1.24 per SD), \textbf{mediating 11.74\% (7.91--18.95\%)} of the
score--PH association. The mediation estimate is the part to hold loosely: metabolites were
measured once at baseline, so the mediator is cross-sectional with respect to a 13-year
outcome, and the causal ordering of exposure, metabolome and disease cannot be established
from that design. Residential exposure models add the usual non-differential error, which
biases the 11.74\% downward.}
{10.1016/j.ecoenv.2026.120778}

\band{Neuro, cognition \& mental health\hfill 2 records}{Deep}

\paperentry{Sex modification dementia meta 2026}
{126 million participants: sex does not modify the PM$_{2.5}$--Alzheimer's association}
{Environ Pollut}
{Systematic review and meta-analysis of 24 longitudinal cohorts, \textbf{over 126 million
participants}, searched to September 2025, using the female-to-male ratio of risk ratios as
the modification metric --- the correct estimand, and one that many sex-difference papers
substitute a pair of stratified estimates for. For Alzheimer's disease the pooled
PM$_{2.5}$ \textbf{RRR is 0.963 (95\% CI 0.919--1.009)} and for vascular dementia
\textbf{1.003 (0.954--1.054)}: null, and precisely so. The one positive is all-cause
dementia (1.046, 1.024--1.069), a 4.6\% relative excess in women that the authors decline to
build policy on. Heterogeneity was negligible to low except PM$_{10}$ in Alzheimer's
($I^2=61\%$), no publication bias was detected, and GRADE certainty is moderate. The stated
limitation is the one that matters: methodological heterogeneity in how the constituent
cohorts defined and adjusted sex could mask a true difference, so this is evidence against
current \emph{claims}, not proof of no effect.}
{10.1016/j.envpol.2026.129138}

\paperentry{Passive smoking acute stress 2026}
{Self-reported secondhand smoke and acute stress in 4{,}460 rescue personnel}
{Front Public Health}
{Cross-sectional survey of 4{,}460 grassroots rescue officers and soldiers, screening with
the Acute Stress Disorder scale and the PCL-5. Univariate odds were \textbf{2.509
(1.621--3.883)} for acute stress disorder and 3.298 (1.459--7.451) for PTSD; on full
adjustment the ASD association survived and \textbf{the PTSD association did not}
(2.513, 0.900--5.151). Among never-smokers, exclusive passive exposure still carried
OR 1.910 (1.071--3.408). Three problems limit what can be drawn from this. Exposure is
self-reported with \emph{no particle measurement of any kind}, so the dose is unknown and
likely differentially misreported by people who are already distressed. The design cannot
order exposure and outcome. And the reported multivariate ASD interval --- 1.959--3.144
around a point estimate of 1.990 --- is internally inconsistent, so the univariate estimate
is the one plotted here. Filed as a signal that SHS exposure belongs in trauma screening,
not as a dose-response result.}
{10.3389/fpubh.2026.1790527}

\band{Respiratory \& allergic\hfill 2 records}{Sage}

\paperentry{MobiliSense ambulatory spirometry 2026}
{Sub-hourly personal black carbon and a measurable drop in FEV$_1$}
{Research Square}
{\textbf{Preprint.} 199 adults in metropolitan Paris, 2018--2020, each carrying two
portable pollutant monitors, a GPS receiver and an accelerometer through ordinary daily
mobility, with \textbf{2{,}504 spirometry tests} taken morning and evening over three days.
Mixed-effects models tested lags from 15 minutes to 6 hours. In multipollutant models, a
1\,$\mu$g/m$^3$ increase in BC one hour before testing was associated with FEV$_1$ lower by
\textbf{0.016\,L (95\% CI $-$0.024 to $-$0.008)}, and at two hours by
\textbf{0.021\,L ($-$0.034 to $-$0.007)}; PM$_{2.5}$ at 15, 30 and 60 minutes reduced
FEV$_1$/FVC by 0.60, 0.70 and 0.50 percentage points. O$_3$ ran positive, which is the
expected artefact of O$_3$ anticorrelating with traffic. The threats are real: 199
participants, three days, no personal-monitor validation reported against reference
instruments, and spirometry quality in the field is hard. But this is the design that
resolves an exposure window a fixed monitor cannot see at all.}
{10.21203/rs.3.rs-10386483/v1}

\paperentry{Cardiorespiratory review 2026}
{Narrative review of air pollution and cardiorespiratory disease}
{Med Gas Res}
{Orientation-level review of the mechanisms linking ambient air pollution to
cardiometabolic risk and cardiovascular death, with sections on heart failure, COPD and
asthma exacerbation, on vulnerable groups (pregnancy, neonates, children, pre-existing
cardiorespiratory disease), on the COVID-19 lockdown natural experiment, on international
divergence in air quality standards, and on mitigation from policy through technology.
There is no systematic search protocol, no PROSPERO registration, no pooled estimate and
no risk-of-bias assessment, so it cannot be cited as evidence for any magnitude --- and it
does not claim to be. Listed because the lockdown-period material and the standards
comparison are a useful map of where the regulatory targets currently sit, and because
naming it here prevents it being mistaken for a systematic review later.}
{10.4103/mgr.MEDGASRES-D-25-00201}

\clearpage

\band{Reproductive \& developmental\hfill 2 records}{Sky}

\paperentry{Prenatal greenspace ASD window 2026}
{A mid-to-late pregnancy window for green space --- and PM$_{2.5}$ is not the mechanism}
{Environ Pollut}
{1:3 age- and sex-matched case-control from a tertiary maternal and child health hospital in
central China: 207 children with autism spectrum disorder, 621 controls. Monthly residential
NDVI and PM$_{2.5}$ were estimated from 10 months before to 3 months after birth and entered
a distributed lag nonlinear model, with product terms testing NDVI $\times$ PM$_{2.5}$
interaction, bootstrapped mediation, and Bonferroni correction. Cases had lower NDVI in
\emph{all 14 months}; the protective window localises to mid-to-late pregnancy, roughly 6 to
1 month before term birth. The negative result is the one that earns the entry:
\textbf{PM$_{2.5}$ neither mediated nor modified the green-space effect} after correction,
which closes the most obvious alternative explanation rather than leaving it hanging.
Caveats are the standard ones for this design --- residential address as the exposure proxy,
hospital-based case ascertainment, 207 cases against a 14-month lag structure, and green
space as a proxy for an unmeasured bundle (noise, activity, socioeconomic position).}
{10.1016/j.envpol.2026.129118}

\paperentry{Preeclampsia nationwide 2018-23 2026}
{\textsc{metadata only} --- nationwide air pollution and preeclampsia, 2018--2023}
{Eur J Intern Med}
{\textbf{The PubMed record carries no abstract, and neither Crossref nor Europe PMC held
one on the entry date.} What the record establishes: a nationwide study covering 2018--2023
has examined the association between ambient air pollution and preeclampsia. What it does
not establish: the country, the pollutant set, whether PM was resolved at all, the case
count, the exposure-assignment method, the trimester windows examined, or the direction and
magnitude of any association --- none of which follows from the title, and none of which is
asserted here. Flagged for retrieval. Hypertensive disorders of pregnancy are one of the
better-powered PM endpoints and a nationwide six-year series would be a substantial
addition if the exposure model holds up.}
{10.1016/j.ejim.2026.107175}

\band{Occupational \& indoor exposure\hfill 3 records}{Clay}

\paperentry{PSH smoke-free home RCT 2026}
{A smoke-free-home trial that missed its primary outcome, reported plainly}
{JAMA Netw Open}
{Cluster-randomised trial across 40 permanent-supportive-housing sites in the San Francisco
Bay Area, 400 residents (mean age 54.5, 63\% male) who smoked at home, January 2022 to March
2025, intervention sites receiving one-to-one coaching on adopting a smoke-free home and
staff at all sites receiving brief cessation training. \textbf{The primary outcome failed}:
90-day smoke-free-home adoption was 6.8\% versus 4.8\%, \textbf{OR 1.45 (0.69--3.07)}. What
moved was the softer secondary --- any adoption of at least one day, 63.4\% versus 36.8\%,
\textbf{adjusted OR 3.83 (2.63--5.57)} --- and CO-verified 7-day abstinence, 6.3\% versus
1.0\%, OR 6.94 (1.69--28.45), which rests on \emph{14 events in total} and should not be
read as an effect size. The honest reading is that brief coaching moves intention and not
maintenance in a population with very high baseline exposure, and that a waiting-list
control cannot rule out attention effects on the self-reported outcomes.}
{10.1001/jamanetworkopen.2026.32697}

\paperentry{Dutch school air cleaner RCT 2026}
{Protocol: 180 classrooms, HEPA versus ionisation versus nothing}
{Front Public Health}
{\textbf{Protocol only, no results} (ClinicalTrials.gov NCT07479420). 180 classrooms in 29
Dutch primary schools, clustered within schools and randomised to HEPA-filter portable air
cleaners, ionisation/plasma cleaners, or none. Units were laboratory pre-tested, safety
screened and operated at comparable clean air delivery rates --- which is the design choice
that makes the comparison mean something, since CADR is the confounder that sinks most
field air-cleaner studies. The schedule is a three-week baseline or two-week post-intervention
control plus three repeated three-week intervention periods, to 14 weeks per school. Airborne
dust is collected passively on electrostatic dustfall collectors and assayed for bacterial and
viral markers, with active sampling in 12 classrooms to validate the passive method;
PM$_{10}$, PM$_4$, PM$_{2.5}$, PM$_1$, CO$_2$, temperature, humidity and VOCs are logged
continuously. Absenteeism and parent-reported symptoms are retrospective, which is the
weakest element. Worth tracking to results --- ionisation devices rarely get a fair,
CADR-matched test.}
{10.3389/fpubh.2026.1892902}

\paperentry{CFRP inhalation risk review 2026}
{Carbon-fibre composite dust: fibre-like fragments, ultrafines, and thin evidence}
{Regul Toxicol Pharmacol}
{Life-cycle review of inhalation risk from carbon fibre-reinforced plastics across
manufacturing, machining, recycling and end-of-life processing. Mechanical processing
generates respirable high-aspect-ratio fragments; thermal processing generates ultrafine
particles together with hazardous gases. The review's central judgement is a refusal:
available evidence \textbf{does not demonstrate a magnitude or pattern of inhalation
toxicity comparable to asbestos or to pathogenic high-aspect-ratio carbon nanotubes}, but it
is too limited and heterogeneous --- particularly for chronic exposure --- to support any
quantitative relative-toxicity comparison in either direction. It identifies the actual
gaps: real-world occupational exposure data, harmonised emission scenarios, internal dose
and biopersistence. As CFRP volumes rise with wind-blade and aerospace retirement, this is
a source class whose workplace aerosol nobody is routinely monitoring.}
{10.1016/j.yrtph.2026.106243}

\clearpage

\band{Burden, policy \& mitigation\hfill 3 records}{Coral}

\paperentry{Flue-gas condensable PM removal 2026}
{Removing the particulate mass that filterable-PM controls never see}
{J Hazard Mater}
{Condensable particulate matter leaves coal-fired and municipal-waste-incineration stacks as
vapour and forms particles in the atmosphere downstream, which is why filterable-PM control
devices and, in many jurisdictions, the compliance measurement itself miss it entirely. This
bench study passes flue gas through a string-of-beads liquid and reports \textbf{29.56--65.27\%
overall CPM removal at 1.17--8.25\% of the energy} of conventional condensation. The organic
fraction is removed more readily than the inorganic throughout, with the sulfate fraction
best at \textbf{78.34\%}. Mechanistically, condensation dominates at high flue-gas
temperature and dissolution takes over as the gas cools, and a two-term model combining them
reproduces the measured removal. It is bench scale with synthetic conditions --- no
long-duration fouling, no liquid regeneration economics, no full-scale pressure drop --- so
the efficiency range is an upper bound on what a retrofit would deliver.}
{10.1016/j.jhazmat.2026.143505}

\paperentry{Turkiye stroke GBD 2000-2023 2026}
{Ambient PM is the leading environmental risk for stroke death in Türkiye}
{Neuroepidemiology}
{GBD 2023 secondary analysis for Türkiye. Age-standardised stroke incidence fell from 141.47
to 103.05 per 100{,}000 between 2000 and 2023, deaths from 132.61 to 57.03, and DALYs from
2696.30 to 1073.51 --- a genuine and large decline. In 2023, metabolic risks accounted for
38.24\% of stroke deaths (high blood pressure alone 21.29\%) and smoking 12.20\% among
behavioural risks; \textbf{ambient particulate matter led the environmental risks at
11.93\%}, placing it within a percentage point of tobacco. Macroeconomic regressions gave
mixed and partly counterintuitive results --- universal health coverage service capacity
correlated positively with death rates --- which is a familiar ecological artefact of
ascertainment improving with health-system capacity. The deeper limitation for this watch:
GBD's PM exposure surface is itself a model, so the 11.93\% is an attribution built on
modelled exposure, not an independent measurement.}
{10.1159/ned/adwag030}

\paperentry{Intergenerational mortality equity 2026}
{\textsc{metadata only} --- intergenerational equity in acute and chronic mortality impacts}
{Environ Int}
{\textbf{No abstract retrievable on the entry date} from Crossref per-DOI, Europe PMC or
PubMed. What the record establishes: mortality impacts of \emph{both} acute and chronic air
pollution exposure have been assessed under an intergenerational health-equity framing ---
an unusual decomposition, since acute and chronic burden estimates are normally produced by
separate literatures with incompatible baselines. What it does not establish: the
populations, the pollutant set, whether PM is resolved, the discounting or
life-table assumptions that an intergenerational comparison entirely depends on, the
geography, or any result. None is asserted here. Flagged for retrieval as the highest-value
of the four metadata-only records in this issue.}
{10.1016/j.envint.2026.110515}

\band{Mechanistic toxicology\hfill 1 record}{Violet}

\paperentry{Tire wear particle phytotoxicity 2026}
{Tire wear particles stunt a roadside grass and raise root oxidative stress}
{Curr Biol}
{Controlled tire-wear-particle exposure of \textit{Bromus catharticus}, a common roadside
weed in Japan, testing an endpoint the TWP literature has largely skipped --- previous work
concentrated on aquatic systems and on crops, not on the roadside and urban-park vegetation
that sits metres from the source. Both \textbf{shoot and root growth were significantly
inhibited}, and root malondialdehyde --- a lipid-peroxidation marker --- rose significantly,
indicating oxidative stress as at least one operative mechanism. The abstract does not give
the applied dose, the particle size distribution, the exposure route (soil versus foliar) or
the exposure duration, and all four determine whether the treatment is environmentally
realistic; a single annual grass in a pot cannot stand for roadside vegetation generally.
Carried here because TWP is the fastest-growing non-exhaust PM source, is poorly sized by
optical monitors, and now has a documented terrestrial effect at the point of emission.}
{10.1016/j.cub.2026.06.087}

\clearpage

% ==================================================================== TABLE
\section*{\sffamily\bfseries\color{Deep}Corpus register}
\vspace{-1.5mm}{\color{Amber}\rule{\linewidth}{1.1pt}}\vspace{2.5mm}

{\fontsize{6.7}{8.4}\selectfont
\setlength{\tabcolsep}{2.4pt}
\begin{longtable}{@{}>{\RaggedRight}p{25mm}>{\RaggedRight}p{41mm}>{\RaggedRight}p{27mm}%
>{\RaggedRight}p{23mm}>{\RaggedRight}p{18mm}>{\RaggedRight}p{15mm}@{}}
\toprule
\textsf{\textbf{First author}} & \textsf{\textbf{Focus}} & \textsf{\textbf{Journal}} &
\textsf{\textbf{Design}} & \textsf{\textbf{Metric}} & \textsf{\textbf{Region}}\\
\midrule
\endfirsthead
\toprule
\textsf{\textbf{First author}} & \textsf{\textbf{Focus}} & \textsf{\textbf{Journal}} &
\textsf{\textbf{Design}} & \textsf{\textbf{Metric}} & \textsf{\textbf{Region}}\\
\midrule
\endhead
\multicolumn{6}{@{}l@{}}{\cellcolor{Mist}\textsf{\textbf{\textcolor{Deep}{Neuro / mental health}}}}\\[0.6mm]
\href{https://doi.org/10.1001/jamanetworkopen.2026.25015}{Dales} & Wildfire-sourced PM$_{2.5}$ and ED visits for migraine and other primary headache & \textit{JAMA Netw Open} & Case-crossover & PM$_{2.5}$ (wildfire) & Canada\\
\href{https://doi.org/10.1093/ageing/afag215}{Glover} & Umbrella meta-review of urban environmental exposures and cognitive health & \textit{Age Ageing} & Umbrella review / meta-analysis & PM$_{2.5}$ & Global\\
\href{https://doi.org/10.1016/j.nsa.2026.107019}{Goossens} & In utero air pollution exposure and pre-/postnatal brain development & \textit{Neurosci Appl} & Narrative review & PM$_{2.5}$ / PM$_{10}$ / PM$_{\mathrm{coarse}}$ & Global\\
\href{https://doi.org/10.1016/j.isci.2026.116617}{Ma} & PM$_{2.5}$-induced neurotoxicity and depression: exposure risks and mechanisms & \textit{iScience} & Narrative review & PM$_{2.5}$ + components & Global\\
\href{https://doi.org/10.3760/cma.j.cn112150-20250804-00752}{Pan} & Maternal PM exposure and offspring neurodevelopment: effects and mechanisms & \textit{Zhonghua Yu Fang Yi Xue Za Zhi} & Narrative review & PM (bulk) & China\\
\addlinespace[0.9mm]
\multicolumn{6}{@{}l@{}}{\cellcolor{Mist}\textsf{\textbf{\textcolor{Teal}{Cardiovascular \& metabolic}}}}\\[0.6mm]
\href{https://doi.org/10.3760/cma.j.cn112338-20251120-00843}{Huang} & PM$_{2.5}$ and dyslipidaemia in rural Chinese children; dietary effect modification & \textit{Zhonghua Liu Xing Bing Xue Za Zhi} & Cross-sectional (multistage) & PM$_{2.5}$ & China\\
\href{https://doi.org/10.32394/pe/218966}{Kazmierska} & Air pollution and cardiovascular disease: review of current data & \textit{Przegl Epidemiol} & Narrative review & PM$_{2.5}$ / PM$_{10}$ & Global\\
\href{https://doi.org/10.1007/s00125-026-06814-2}{Xi} & Air pollution, built environment and temperature vs. prediabetes phenotypes & \textit{Diabetologia} & Prospective cohort & PM$_{2.5}$ & Germany\\
\href{https://doi.org/10.1161/JAHA.125.046917}{Yang} & Short-term PM$_{2.5}$ and hospitalisation/ED use in atrial fibrillation & \textit{J Am Heart Assoc} & Case-crossover & PM$_{2.5}$ (wildfire-influenced) & USA\\
\addlinespace[0.9mm]
\multicolumn{6}{@{}l@{}}{\cellcolor{Mist}\textsf{\textbf{\textcolor{Sky}{Reproductive \& developmental}}}}\\[0.6mm]
\href{https://doi.org/10.1016/j.jes.2025.10.043}{Chen} & Prenatal air pollution, fetal growth trajectory phenotypes and child neurodevelopment & \textit{J Environ Sci (China)} & Prospective cohort & PM$_{2.5}$ / PM$_{10}$ & China\\
\href{https://doi.org/10.1016/j.ecoenv.2026.120566}{Huang} & Joint prenatal bisphenol + air-pollutant mixture exposure and preschool IQ & \textit{Ecotoxicol Environ Saf} & Prospective cohort (mixtures) & PM$_{2.5}$ / PM$_{10}$ & China\\
\href{https://doi.org/10.1007/s40572-026-00552-8}{Peterson} & Ambient air pollution, PAHs and pubertal development: critical review & \textit{Curr Environ Health Rep} & Critical review & PM$_{2.5}$ components & Global\\
\href{https://doi.org/10.1016/j.envres.2026.125313}{Xie} & Long-term PM$_{2.5}$ and timing of pubertal development in boys and girls & \textit{Environ Res} & Prospective cohort & PM$_{2.5}$ & China\\
\addlinespace[0.9mm]
\multicolumn{6}{@{}l@{}}{\cellcolor{Mist}\textsf{\textbf{\textcolor{Sage}{Respiratory \& allergic}}}}\\[0.6mm]
\href{https://doi.org/10.1007/s00484-026-03289-z}{Bayar Muluk} & Climate change and otolaryngology: airway disease impacts and emission reduction & \textit{Int J Biometeorol} & Narrative review & PM$_{2.5}$ & Global\\
\href{https://doi.org/10.4103/mgr.MEDGASRES-D-25-00164}{Liao} & Short-term air pollution, respiratory symptoms and lung function in asthma & \textit{Med Gas Res} & Panel study & PM$_{2.5}$ / PM$_{10}$ & China\\
\href{https://doi.org/10.3390/epidemiologia7040097}{Sanchez de Prada} & Ambient PM and asthma-coded admissions in children 0-3 y, Spain & \textit{Epidemiologia} & Ecological panel & PM$_{2.5}$ / PM$_{10}$ & Spain\\
\href{https://doi.org/10.1016/j.jhazmat.2026.143076}{Yang} & Allergenic vs. toxic risk of airborne seed hairs compared with ambient PM$_{2.5}$ & \textit{J Hazard Mater} & Field measurement & PM$_{2.5}$ + bioaerosol & China\\
\addlinespace[0.9mm]
\multicolumn{6}{@{}l@{}}{\cellcolor{Mist}\textsf{\textbf{\textcolor{Violet}{Mechanistic toxicology}}}}\\[0.6mm]
\href{https://doi.org/10.1016/j.jes.2025.12.074}{Ding} & PM$_{2.5}$ promotes vascular calcification via miR-27b-5p/HES1 signalling & \textit{J Environ Sci (China)} & In vitro + organoid + murine & PM$_{2.5}$ & China\\
\href{https://doi.org/10.1007/s00360-026-01698-5}{Maraschi} & Osmoregulatory disruption and metal bioaccumulation from settleable atmospheric PM & \textit{J Comp Physiol B} & Controlled exposure (ecotox) & Settleable APM & Brazil\\
\href{https://doi.org/10.1016/j.jes.2025.11.018}{Wu} & Macrophage NLRP3 inflammasome in real-ambient PM-induced hepatic glucose dysregulation & \textit{J Environ Sci (China)} & Real-ambient murine + co-culture & Real-ambient PM & China\\
\addlinespace[0.9mm]
\multicolumn{6}{@{}l@{}}{\cellcolor{Mist}\textsf{\textbf{\textcolor{Clay}{Occupational \& indoor}}}}\\[0.6mm]
\href{https://doi.org/10.3389/fmedt.2026.1796828}{Chen} & Nested suction instrument to cut intraoperative surgical-smoke PM$_{2.5}$ & \textit{Front Med Technol} & Bench / device experiment & PM$_{2.5}$ (surgical smoke) & China\\
\href{https://doi.org/10.1093/annweh/wxag055}{Cole} & Determinants of respirable crystalline silica exposure in Australian tunnelling & \textit{Ann Work Expo Health} & Observational exposure study & Respirable dust / RCS & Australia\\
\href{https://doi.org/10.1016/j.jhazmat.2026.143081}{Ding} & Stage-resolved hair Zn-isotope dynamics as a firefighting combustion-exposure tracer & \textit{J Hazard Mater} & Repeated-measures biomarker & Combustion-derived PM & China\\
\href{https://doi.org/10.3390/jox16040119}{Ramsden} & Xenobiotic hazards in aircraft cabin air & \textit{J Xenobiot} & Review + risk economics & Ultrafine particles & Global\\
\addlinespace[0.9mm]
\multicolumn{6}{@{}l@{}}{\cellcolor{Mist}\textsf{\textbf{\textcolor{Amber}{Sensing, forecasting \& instrumentation}}}}\\[0.6mm]
\href{https://doi.org/10.1016/j.jenvman.2026.130515}{Shi} & Adaptive spatiotemporal graph transformer for multi-site multi-step PM$_{2.5}$ forecasting & \textit{J Environ Manage} & Deep-learning model & PM$_{2.5}$ & China / India\\
\href{https://doi.org/10.3390/bios16070380}{Tsuchiya} & Seasonal and indoor-outdoor characteristics of fluorescent aerosol particles in offices & \textit{Biosensors} & Real-time sensor campaign & Bioaerosol / coarse & Japan\\
\addlinespace[0.9mm]
\multicolumn{6}{@{}l@{}}{\cellcolor{Mist}\textsf{\textbf{\textcolor{Amber}{Exposure assessment \& modelling}}}}\\[0.6mm]
\href{https://doi.org/10.1029/2025GH001621}{Solorzano-Araque} & Multiscale exposure modelling: LOTOS-EUROS + mobility surveys + particle toxicity & \textit{GeoHealth} & CTM + mobility integration & PM$_{2.5}$ / PM$_{10}$ & Colombia\\
\addlinespace[0.9mm]
\multicolumn{6}{@{}l@{}}{\cellcolor{Mist}\textsf{\textbf{\textcolor{Coral}{Burden, policy \& mitigation}}}}\\[0.6mm]
\href{https://doi.org/10.1002/hsr2.72886}{Karamzad} & Ischaemic heart disease burden in MENA 1990-2021; PM attributable fraction 28.2\% & \textit{Health Sci Rep} & GBD secondary analysis & Ambient PM & MENA\\
\href{https://doi.org/10.1371/journal.pone.0347989}{Rees} & HEPA filtration to reduce respiratory infection in care homes & \textit{PLoS One} & RCT process evaluation & Filtered PM & UK\\
\href{https://doi.org/10.1016/j.jes.2026.01.009}{Shan} & Policy-driven emission reductions and post-pandemic rebound, high-resolution inventories & \textit{J Environ Sci (China)} & Emission inventory & PM emissions & China\\
\addlinespace[0.9mm]
\multicolumn{6}{@{}l@{}}{\cellcolor{Mist}\textsf{\textbf{\textcolor{Slate}{Other clinical endpoints}}}}\\[0.6mm]
\href{https://doi.org/10.3390/dj14070440}{Lupita} & PM$_{2.5}$, preventive-care capacity and caries/edentulism burden across the EU & \textit{Dent J} & Ecological panel & PM$_{2.5}$ & EU-27\\
\href{https://doi.org/10.1016/j.jdin.2026.06.008}{Park} & Long-term PM exposure and risk of androgenetic alopecia & \textit{JAAD Int} & Nationwide cohort & PM$_{2.5}$ / PM$_{10}$ & South Korea\\
\href{https://doi.org/10.3171/2026.3.PEDS2615}{Sun} & Environmental exposure and area deprivation vs. paediatric hydrocephalus outcomes & \textit{J Neurosurg Pediatr} & Registry cohort + geospatial & PM$_{2.5}$ (satellite-derived) & USA\\
\addlinespace[0.9mm]
\bottomrule
\end{longtable}}


\vspace{4mm}
\begin{keybox}
{\sffamily\bfseries\footnotesize\color{Deep}Provenance and method}\\[1.2mm]
{\fontsize{7.2}{8.8}\selectfont
\textbf{Window.} Entry date 9 September 2026 only. The 9 September scheduled run did not
fire; this issue was built on 10 September at 18:39 local, which is before the 22:00 gate,
so 9 September is the newest buildable date and 10 September is deliberately left
unopened rather than harvested against a partial day.\par
\vspace{0.8mm}
\textbf{Sources.} \textbf{PubMed} E-utilities on \texttt{[Date - Entry]}, two axes run as
separate queries (health 11 hits, sensing 1); the \textbf{PubMed connector} run
independently on the same two axes returned 16 on the health axis, \textbf{nine PMIDs the
E-utilities leg never returned}, three of which ship --- the first run in which the
connector contributed net-new records rather than confirming the harvester.
\textbf{Europe PMC} \texttt{CREATION\_DATE} (24 hits, including two preprints that ship).
\textbf{Crossref by ISSN} across the watch's journal list (68 hits) --- the only leg that
reaches \textit{Atmos.\ Meas.\ Tech.}, \textit{Atmos.\ Chem.\ Phys.} and
\textit{Atmosphere}, which PubMed does not index and which supplied four of the seven
sensing records. \textbf{arXiv} relevance query with a local recency screen --- one record
passed it (wildfire smoke and Californian crop yields) and proved to be already in the
archive by title. \textbf{OpenAlex} returned HTTP 429 and contributed nothing; that leg has been
degraded since July. \textbf{Consensus} returned three sensing hits: one is already in
\texttt{seen.json} from 5 August, and the other two were checked against their Crossref
\emph{created} dates --- 26 March and 20 January 2026 --- so both sit far outside this
window and are archive gaps rather than records for this issue.\par
\vspace{0.8mm}
\textbf{Screening.} 105 raw records, 102 after cross-source deduplication on
PMID\,$\rightarrow$\,DOI\,$\rightarrow$\,normalised title, 14 already in
\texttt{seen.json}, 88 fresh, plus 9 connector-only records.
\textbf{26 in scope, 71 excluded}, each with a logged reason in
\texttt{state/rejected.jsonl}. Nineteen exclusions are \textit{Meas.\ Sci.\ Technol.}
bycatch from the ISSN leg (bearing diagnostics, SLAM, inertial navigation) and eleven are
non-atmospheric \textit{ES\&T}.\par
\vspace{0.8mm}
\textbf{Metadata-only records.} Four Elsevier records had no abstract in existence anywhere
on the entry date --- checked per-DOI against Crossref, against Europe PMC, and against
PubMed \texttt{efetch}. Crossref indexes a DOI before the publisher deposits the abstract,
and a one-day window sits inside that gap; this is the second consecutive issue to hit it
and it should now be treated as expected behaviour for same-day Elsevier deposits, not as
a fault. These four are unambiguously in scope on the title and ship as \textsc{metadata
only}, tier C, asserting no finding. Two further abstract-free records were \emph{rejected}
because without an abstract their PM scope could not be confirmed at all.\par
\vspace{0.8mm}
\textbf{Trial watch.} ClinicalTrials.gov returned \textbf{0} PM or air-pollution
registrations or updates on 9 September against \textbf{1{,}189 registry-wide} on the
no-term control --- a topical zero, not an outage. Accumulated trial watch surfaces in the
weekly rollup (next: Saturday 12 September).\par
\vspace{0.8mm}
\textbf{Labels and estimates.} Subtopic, design, particle-metric and geography labels are
assigned from title, abstract, keywords and MeSH; assignment is single-label by dominant
contribution and is a judgement, not a controlled vocabulary. Effect estimates are
transcribed verbatim from abstracts and use heterogeneous contrasts --- see the forest-plot
caption, which also records the one estimate omitted and why. All DOIs resolve to the
publisher of record and are checked automatically before build. Abstracts remain the
copyright of their publishers.\par}
\end{keybox}

\end{document}
